\documentclass[preprint,12pt]{elsarticle}
\usepackage[T1]{fontenc}\usepackage{lmodern}\usepackage{booktabs}\usepackage{amsmath}
\renewcommand{\thetable}{S\arabic{table}}
\begin{document}
\begin{center}{\Large Supplementary material}\\[4pt]
Eight-bit quantization of a tiny wrist-inertial network for Parkinson's disease detection: calibration sensitivity, evaluation pitfalls and microcontroller validation\end{center}

\section*{S1. Runs affected by a script error}
In five runs, the batch normalization after each depthwise convolution had been disabled by a script error (21,841 instead of 22,737 parameters). These runs were superseded by the runs of the specified network in the main text and are listed here for completeness (Table~S1). In these runs, PTQ lost more than 0.05 AUROC in 9 of 25 fold-run units and QAT in none.

\begin{table}[h]
\caption{Within-fold clean AUROC of FP32 and mean difference from FP32 (worst fold) for the runs without BN after the depthwise convolutions, PD versus HC.}
\centering\scriptsize\setlength{\tabcolsep}{2pt}
\begin{tabular}{@{}lllccccc@{}}
\toprule
Run & Filter & Seed & FP32 & PTQ int8 & PTQ, 2000 & PTQ $16\times8$ & QAT int8 \\
\midrule
s1 & Zero-phase & 1 & 0.862 & $-0.060$ ($-0.095$) & -- & -- & $+0.005$ ($-0.008$) \\
s2 & Zero-phase & 2 & 0.881 & $-0.050$ ($-0.124$) & -- & -- & $-0.001$ ($-0.050$) \\
c0 & Causal & 0 & 0.871 & $-0.030$ ($-0.119$) & $-0.021$ & $-0.001$ & $+0.021$ ($-0.006$) \\
c1 & Causal & 1 & 0.866 & $-0.007$ ($-0.020$) & $-0.015$ & $-0.003$ & $+0.024$ ($+0.002$) \\
c2 & Causal & 2 & 0.861 & $-0.037$ ($-0.068$) & $-0.016$ & $0.000$ & $+0.010$ ($-0.005$) \\
\bottomrule
\end{tabular}
\end{table}

\section*{S2. Exploratory external-cohort analysis}
Added after the first results and run only with the script error of Section~S1, this analysis trained an accelerometer-only network (20,497 parameters) on all PADS participants and tested it on the Monipar cohort (21 PD, 7 HC). Participant-level AUROC was 0.86 (95\% CI 0.70 to 0.98) for FP32, 0.95 for PTQ and 0.86 for QAT. With seven controls, these intervals are too wide to inform any comparison, and the analysis was not repeated with the specified network.

\section*{S3. Per-window end-to-end firmware}
An earlier firmware ran decimation, causal filtering, channel construction, normalization, quantization and inference for one window at a time with the fold-0 QAT model of run c0 (Section~S1). It took 14.6~ms per window (2.4~ms preprocessing, 12.2~ms inference) with 42.6~kB of RAM. Its inference time differs from the 10.7~ms of the whole-fold replay firmware, whose build and model differ; the difference was not decomposed. The streaming firmware of the main text, run with the specified causal model k0, supersedes this test.

\section*{S4. Range and fidelity measurements on retrained fold-0 models}
The range measurements on the ten retrained fold-0 models of runs q0 and q1 (FP32 AUROC 0.830 to 0.919) showed that the loss is not caused by input saturation. With default min/max calibration, at most 0.005\% of test input values saturated, and across draws the loss was not correlated with the input scale (Spearman $\rho=-0.03$ for 300 windows and $+0.14$ for 2000 windows) or with the output scale. Over 100 draws per setting, the loss averaged 0.023 (300 windows) and 0.051 (2000 windows), exceeded 0.05 in 15\% and 41\% of draws and reached 0.207 and 0.317. Clipping the same draws at the 0.1/99.9 input percentiles made the input scale three to five times finer (0.074 versus 0.224 and 0.338), saturated 0.07\% of test input values, and reduced the mean loss to 0.010 and 0.009, the share of draws losing more than 0.05 to 2\% and the worst draw to 0.059 and 0.057. The worst draw of every model was within 0.06 with clipping, against up to 0.207 (300 windows) and 0.317 (2000 windows) without it, and the improvement in the per-model mean loss was clearest for 2000 windows (better in 9 of 10 models, Wilcoxon $p=0.006$; 300 windows, 7 of 10, $p=0.19$). In the same models, $16\times8$ changed AUROC by at most 0.003 and QAT by $-0.022$ to $+0.030$. Some draws gained up to 0.043 AUROC over FP32, and FP32 AUROC itself varied by about 0.04 between retrainings with the same folds; with about 16 controls per test fold, a single-draw difference below about 0.05 is within evaluation noise, which is why we report the share of draws exceeding 0.05. The ten models of runs v0 and v1 (FP32 AUROC 0.783 to 0.892) reproduced the pattern and showed where it acts. With min/max calibration, 17\% (300 windows) and 31\% (2000 windows) of draws lost more than 0.05 participant AUROC, against 0 to 2\% with clipping; window-level AUROC, which rests on about 8400 windows, fell by 0.049 and 0.068 on average with min/max calibration and by 0.012 and 0.013 with clipping at 1/99, and window decisions agreed with FP32 in 79\% and 71\% of windows against 91\% and 90\%. Within a model, draws with a larger logit error lost more participant AUROC (Spearman $\rho=-0.40$) and the window- and participant-level losses were correlated ($\rho=+0.51$), so the participant-level losses are not artefacts of the small number of test controls. Min/max calibration produced activation scales 1.6 to 2.2 times larger than clipped calibration in the convolutional and residual layers, 1.4 times larger after pooling and 1.2 times at the output, and the input scale 5.5 times larger. Within a model, larger scales went with larger losses in every layer ($\rho=-0.10$ to $-0.28$ for min/max draws), most strongly in the stem, the first pointwise convolutions and the first residual additions and least in the output layer. 16$\times$8 changed participant AUROC by at most 0.005, and QAT by $-0.030$ to $+0.039$.
\end{document}
